\begin{textsample}{POS Dim 3 – es – Score 12.00 – es/es\_inf\_9.txt}  \label{ex:f3_pos_070}
Para Refletir A ação docente deve promover práticas que estejam voltadas para a observação e \textbf{escuta} atenta dos interesses, \textbf{desejos} e necessidades das \textbf{crianças}.
Com a intenção de garantir as aprendizagens e o desenvolvimento por meio dos objetivos deste campo de experiências, o trabalho pedagógico da Educação \textbf{Infantil} ganha força ao promover situações que contemplem experiências em relação ao \textbf{cuidado} de si mesmo, ao autoconhecimento e relacionamento com o outro.
Assim, algumas reflexões se fazem necessárias nesse processo : A organização da prática pedagógica tem como objetivo garantir qualidade nas interações dos \textbf{bebês}, das \textbf{crianças} bem \textbf{pequenas} e das \textbf{crianças} \textbf{pequenas}, de modo que, conforme vivem esse processo, nos diferentes tempos e espaços da instituição de Educação \textbf{Infantil}, possam aprender a distinguir e a expressar sensações, percepções, preferências, sentimentos, emoções e pensamentos/opiniões?
Durante o planejamento, os espaços escolares são pensados como ambientes que proporcionam múltiplas experiências, que transmitem segurança e que devem ser aconchegantes e diversos, estimulando o desenvolvimento da autonomia em relação ao \textbf{cuidado} consigo e com o outro?
A \textbf{criança} é \textbf{apoiada} na construção do entendimento pelo \textbf{cuidado} com sua saúde e bem-estar, pela preservação do ambiente escolar, pela criação de \textbf{hábitos} ligados à \textbf{limpeza}, à coleta de lixo produzido na realização das atividades cotidianas e à reciclagem dos inservíveis?
São planejadas situações de interações positivas que ajudem as \textbf{crianças} a criarem, construírem relações de confiança e amizade?
São propostos materiais, atividades e \textbf{brincadeiras} em que as \textbf{crianças} percebam a necessidade de compartilhar, cooperar e socializar, de forma a ajudá -las a reconhecer a existência do ponto de vista do outro, considerando opiniões, sentimentos e intenções do outro, construindo atitudes negociadoras e tolerantes?
São considerados os momentos de adaptação, seja no início de ano letivo ou de \textbf{crianças} que chegam no decorrer do ano letivo, com situações tranquilas e que contribuam com a criação de vínculos entre as \textbf{crianças}?
É importante organizar o ambiente e as \textbf{rotinas} com intencionalidade pedagógica ajudando a \textbf{criança} a desenvolver o sentimento de autoestima, confiança em suas potencialidades, pertencimento a um grupo étnico-racial, identidade pessoal, crença religiosa, local de nascimento, assim como fortalecer os vínculos afetivos com os familiares, oportunizando o acesso a diferentes tradições culturais para compreensão de si mesma e do mundo.

% matched lemmas: apoiar, bebê, brincadeira, criança, cuidado, desejo, escuta, hábito, infantil, limpeza, pequeno, rotina
\end{textsample}
